\documentclass[a4paper,11pt]{article}

\usepackage[utf8]{inputenc}
\usepackage{amsmath,amssymb}
\usepackage[a4paper,margin=2.5cm]{geometry}
\usepackage{enumitem}

\title{Standing Waves and Superposition}
\author{}
\date{}

\begin{document}

\maketitle

\section{Starting point: large rope}

\section{When Waves Meet}

\subsection*{Guiding Question}


textbf{Guiding Question}

Yesterday we learned that sound consists of travelling pressure and density waves.

When you pluck a guitar string, you can hear a clear note for several seconds.

\textbf{Discussion}
\begin{itemize}
    \item If sound waves simply travel away, why does the note last for several seconds?
    \item What must happen inside the instrument while we continue to hear the sound?
    \item Can a single travelling wave explain this observation?
\end{itemize} 

\textbf{Expected ideas}
\begin{itemize}
    \item The string (or the air inside the instrument) must continue vibrating.
    \item The vibration cannot simply leave the instrument immediately.
    \item There must be a wave pattern that remains inside the instrument while sound is continuously radiated into the surrounding air.
\end{itemize}


\textbf{Conclusion}  Something must happen that keeps the vibration inside the instrument instead of letting it simply travel away.
A travelling wave transports the disturbance through the medium.It does not remain in one place. Therefore, a different type of wave is needed to explain how a musical instrument can produce a long, steady note.

\subsection*{Activity 1: A Travelling Wave}

A long rope is stretched along the room, a single pulse is generated. 

\textbf{Discussion}
\begin{itemize}
    \item What happens to the pulse?
    \item Does it stay in one place?
    \item What happens after it reaches the end of the rope?
\end{itemize}


\textbf{Expected answers}
\begin{itemize}
    \item The pulse travels along the rope.
    \item It does not stay in one place.
    \item Each point of the rope moves only for a short time before returning to rest.
    \item When the pulse reaches the end, it disappears (or is reflected, depending on the setup).
    \item The disturbance moves through the rope, not the rope itself.
\end{itemize}

Now, a harmonic wave is generated. 
\textbf{Discussion}
\begin{itemize}
    \item What is different from a single pulse?
    \item Does the wave remain on the rope?
    \item Could this kind of wave produce a steady musical note? Why or why not?
\end{itemize}


\textbf{Expected answers}
\begin{itemize}
    \item The pulse travels along the rope.
    \item It does not stay in one place.
    \item Each point of the rope moves only for a short time before returning to rest.
    \item When the pulse reaches the end, it disappears (or is reflected, depending on the setup).
    \item The disturbance moves through the rope, not the rope itself.
\end{itemize}

Students will notice that the wave simply travels along the rope and eventually disappears at the end.

\vspace{0.5em}
\subsection*{Activity 2: Reflection}

Generate a single pulse with reflection at the end 

\textbf{Discussion}
\begin{itemize}
    \item What happens when the pulse reaches the end?
    \item Does it disappear completely?
    \item Where does the reflected pulse travel?
\end{itemize}

\textbf{Expected answers}
\begin{itemize}
    \item The pulse is reflected.
    \item It travels back along the rope.
    \item For a short time, there are two pulses travelling in opposite directions.
\end{itemize}

Sketch on the board: 

$
\rightarrow \qquad\qquad \leftarrow
$
Question: 

\begin{itemize}
    \item How many waves are now travelling on the rope?
\end{itemize}

Expected answer:

\begin{itemize}
    \item Two waves travelling in opposite directions.
\end{itemize}


\subsection*{Activity 3: Superposition in the Ripple Tank}

Start with a  \textbf{single wave source}, just for a few seconds. 

\textbf{observation:}
\begin{itemize}
    \item transversal wave 
    \item spherical direction of travel 
    \item in the ripple tank: water can spread in two dimensions; can move forward, backward and left right 
    \item rope behaves like a line out of the ripple tank 
    \item only difference: wave can travel in two directions in water but only in one direction in the rope
\end{itemize}

If possible: plain waves which are reflected by a straight barrier otherwise switch on two point sources. 

\textbf{Discussion}
\begin{itemize}
    \item What do you notice?
    \item Are all parts of the water moving equally?
    \item Why are some regions almost calm?
    \item Why are other regions moving much more strongly?
\end{itemize}

\textbf{Expected answers}
\begin{itemize}
    \item Two waves can exist at the same place.
    \item Sometimes the waves reinforce each other.
    \item Sometimes they cancel each other.
    \item The total displacement is the sum of both waves.
\end{itemize}

Introduce the principle of \textbf{superposition}.

\begin{quote}
Whenever two waves meet, their displacements simply add together.
\end{quote}

\section{back to rope}: 
We can also see this phenomenon with ropes $\Rightarrow$ standing waves: 

\textbf{Discussion}

\begin{itemize}
    \item What do you notice?
    \item Which points never move?
    \item Which points move the most?
    \item (Can you explain these observations using the idea of superposition?)
\end{itemize}

\textbf{Expected answers}

\begin{itemize}
    \item Some points always cancel because the two waves have opposite displacements.
    \item These points are called \textbf{nodes}.
    \item At other points the two waves always reinforce each other.
    \item These points are called \textbf{antinodes}.
\end{itemize}

\textbf{Conclusion}

A standing wave is not a new type of wave.

It is formed when two identical harmonic travelling waves move in opposite directions and continuously superpose.

The incoming wave and the reflected wave together produce the standing wave pattern.

\section{Do standing waves occur for all frequencies/wavelengths??}

Form groups, each group gets a rope. 
Tasks: 
\begin{itemize}
    \item Find as many stable wave patterns as possible.
    \item Draw them.
    \item Count the antinodes.
    \item  optional: Measure the wavelength.
    \item optional: Which frequencies work?
    \item optional: Which frequencies do not work?
    \item Can you find a rule?
\end{itemize}

begin{itemize}
    \item Only certain wavelengths produce stable standing waves.
    \item The wavelength must fit exactly onto the string.
    \item As the number of antinodes increases, the wavelength becomes shorter.
\end{itemize}

\textbf{Rule}

A standing wave can only exist if an integer number of half wavelengths fits exactly into the length of the string. (\textbf{je nach festem/losem Ende anpassen })

\section{Seilwellenmaschine}
\section{guitar string}

\begin{itemize}
    \item The fundamental mode has the largest wavelength.
    \item Higher modes have shorter wavelengths.
    \item Since the wave speed is constant,
    $
        c=f\lambda,
    $
    a shorter wavelength means a higher frequency.
    \item short strings produce higher frequencies 
\end{itemize}

\textbf{Conclusion}

Different standing wave patterns correspond to different wavelengths.
Because the wave speed remains constant, each wavelength corresponds to one particular frequency.

These frequencies are the musical notes that we hear.


\section{second day}
\subsection*{Guiding Question}

We have discovered that a stretched string can only support certain standing waves.

\textbf{Discussion}

\begin{itemize}
    \item Do you think standing waves only exist on strings?
    \item What other parts of a musical instrument could support standing waves?
\end{itemize}

\textbf{Expected answers}

\begin{itemize}
    \item Air columns.
    \item The body of an instrument.
    \item Any medium in which waves can travel.
\end{itemize}

\textbf{not all musical instruments use strings, some musical instruments also use vibrating air columns}

\begin{itemize}
    \item string only allows certain standing waves
    \item both ends are fixed $\rightarrow $ no oscillation possible  $\rightarrow $  nodes
    \item recorder: Does it have two fixed ends? - no
    \item depending on the instrument, we have different combinations of fixed/open ends  - called boundary conditions
    \item two fixed ends: guitar 
    \item two open end: transverse flute 
    \item one open and one fixed end: pan flute
\end{itemize}

\textbf{different instruments have different ends and support different wavelength/ frequencies }

\subsection*{Group Investigation}

Divide the class into three groups.

Each group investigates a different type of musical instrument.

\vspace{0.5em}

\textbf{Group 1: String (fixed -- fixed)}

Examples:
\begin{itemize}
    \item Guitar
    \item Violin
    \item Piano
\end{itemize}

\textbf{Tasks}

\begin{itemize}
    \item Draw the first three standing waves.
    \item Mark the nodes and antinodes.
    \item Which wavelengths fit onto the string?
    \item Which frequencies are possible?
\end{itemize}

\vspace{0.5em}
textbf{Group 2: Open Pipe (open -- open)}

Examples:
\begin{itemize}
    \item Flute
    \item Open organ pipe
\end{itemize}

\textbf{Tasks}

\begin{itemize}
    \item Draw the first three standing waves.
    \item Mark the nodes and antinodes.
    \item Which wavelengths fit into the pipe?
    \item Which frequencies are possible?
\end{itemize}

\vspace{0.5em}

\textbf{Group 3: Closed Pipe (closed -- open)}

Examples:
\begin{itemize}
    \item Bottle
    \item Clarinet (approximately)
    \item Closed organ pipe
\end{itemize}

\textbf{Tasks}

\begin{itemize}
    \item Draw the first three standing waves.
    \item Mark the nodes and antinodes.
    \item Which wavelengths fit into the pipe?
    \item Which frequencies are possible?
    \item What is different from the open pipe?
\end{itemize}


\begin{table}[h]
\centering
\renewcommand{\arraystretch}{1.5}
\begin{tabular}{|l|c|c|c|}
\hline
\textbf{Instrument} & \textbf{Ends} & \textbf{Smallest wave} & \textbf{Higher modes} \\
\hline
Guitar string &
Fixed -- Fixed &
Half a wavelength &
1, 2, 3, 4, ...
\\
\hline
Flute &
Open -- Open &
Half a wavelength &
1, 2, 3, 4, ...
\\
\hline
Bottle / Clarinet &
Closed -- Open &
Quarter wavelength &
1, 3, 5, ...
\\
\hline
\end{tabular}
\end{table}


The standing wave with the lowest possible frequency is called the
\textbf{fundamental mode} (or \textbf{fundamental frequency}).

It determines the pitch of the musical note that we hear.

Standing waves with higher frequencies are called
\textbf{harmonics} (or \textbf{overtones}). We will investigate their  significance  in the next days. 


\section{game - optional}
each student is either a node or anti-node. - They have to build standing waves for the  different harmonics and different boundary conditions. 
 
\end{document}
